% P15.28: Durrett's complete density-martingale and product-measure route.
\chapter{Density Martingales and Infinite Product Measures}
\label{sec:p15-density-martingales}

Densities computed from increasing finite information need not converge
in mean. Their lost mass identifies a singular part of the full measure.
For independent observations with equivalent coordinate laws, the product
criterion decides whether two complete sequence laws have the same null
events or are concentrated on disjoint measurable sets. This chapter supplies
Durrett's Section~4.3.3, including its mixture-measure proof, partition
example, tail argument and square-root product criterion
\cite[pp.~227--230]{DurrettPTEAuthor2019}.

\section{Densities under increasing information}

\begin{hbdefinition}{def:p15-density-filtration}{The measures, versions and generated sigma-field}
Let \(\mu\) be a finite nonnegative measure and \(\nu\) a probability
measure on \((\Omega,\mathcal F)\), and put \(c=\mu(\Omega)\).
Write \(\mu\sim\nu\) for mutual absolute continuity.
Let \((\mathcal F_n)_{n\geq0}\) increase with
\(\sigma(\bigcup_n\mathcal F_n)=\mathcal F\). Denote restrictions by
\(\mu_n=\mu|_{\mathcal F_n}\), \(\nu_n=\nu|_{\mathcal F_n}\).
Assume \(\mu_n\ll\nu_n\) for every \(n\).
Choose a finite nonnegative \(\mathcal F_n\)-measurable version
\[
                 X_n=\frac{d\mu_n}{d\nu_n},
                 \qquad X=\limsup_{n\to\infty}X_n\in[0,\infty].
\]
An integrable density admits such a version by changing its infinite
values on a null set. Almost-sure assertions always name their measure.
In particular, changing \(X_n\) on an \(\mathcal F_n\)-measurable
\(\nu\)-null set also changes it on a \(\mu\)-null set, by local absolute
continuity. Countably many choices therefore leave \(X\) unchanged
both \(\mu\)-almost surely and \(\nu\)-almost surely.
\index{density martingale}
\end{hbdefinition}

\begin{hblemma}{thm:p15-density-martingale}{proved=annex:p15.density-martingale}{Restriction densities form a martingale}
Under \cref{def:p15-density-filtration}, \(X_n\) is a nonnegative
martingale under \(\nu\), with \(E_\nu X_n=c\).
It converges to a finite value \(\nu\)-almost surely, so
\[
                       \nu(X=\infty)=0,\qquad E_\nu X\leq c.
\]
This is the complete restriction argument of source Lemma~4.3.6.
\end{hblemma}

\begin{hblemma}{thm:p15-density-mixture-reference}{proved=annex:p15.density-mixture-reference}{A common bounded reference and its limiting densities}
Put
\[
 \rho=\frac{\mu+\nu}{c+1},\qquad \rho_n=\rho|_{\mathcal F_n},
\]
\[
 Y_n=\frac{d\mu_n}{d\rho_n},\qquad
 Z_n=\frac{d\nu_n}{d\rho_n}.
\]
These are nonnegative \(\rho\)-martingales bounded by \(c+1\).
There are finite measurable versions of their limits \(Y,Z\) with,
for every \(A\in\mathcal F\),
\[
 \mu(A)=\int_A Y\,d\rho,\qquad
 \nu(A)=\int_A Z\,d\rho,\qquad
 Y+Z=c+1\quad\rho\text{-almost surely}.
\]
Moreover \(Z_n>0\) \(\rho\)-almost surely and
\(X_n=Y_n/Z_n\) there, simultaneously at all times. Consequently,
\[
 X=Y/Z\text{ on }\{Z>0\},\qquad
 X=\infty\text{ on }\{Z=0\},
 \quad \rho\text{-almost surely}.
\]
The case \(c=0\) is included without dividing by \(c\).
\end{hblemma}

\begin{hbtheorem}{thm:p15-density-lebesgue-decomposition}{proved=annex:p15.density-lebesgue-decomposition}{The singular mass is the infinite-density limit}
For every \(A\in\mathcal F\),
\[
                    \mu(A)=\int_A X\,d\nu+\mu(A\cap\{X=\infty\}).
\]
Thus the two measures
\[
 \mu_r(A)=\int_A X\,d\nu,\qquad
 \mu_s(A)=\mu(A\cap\{X=\infty\})
\]
are the unique Lebesgue decomposition of \(\mu\) with
\(\mu_r\ll\nu\) and \(\mu_s\perp\nu\). In particular
\(X=d\mu_r/d\nu\), \(\nu\)-almost surely. This is source
Theorem~4.3.5, including the measure-specific singular-set assertion
\cite[pp.~227--229]{DurrettPTEAuthor2019}.
\end{hbtheorem}

\begin{hbtheorem}{thm:p15-density-uniform-integrability}{proved=annex:p15.density-uniform-integrability}{Mass preservation and full absolute continuity}
Under the same hypotheses, the following are equivalent:
\[
 \mu\ll\nu;\qquad E_\nu X=c;\qquad
 (X_n)\text{ is uniformly integrable under }\nu;\qquad
 X_n\longrightarrow X\text{ in }L^1(\nu).
\]
In that case \(X=d\mu/d\nu\) and \(X_n=E_\nu(X\mid\mathcal F_n)\).
In general the exact lost mass is
\[
                    \mu_s(\Omega)=c-E_\nu X.
\]
Also \(\mu\perp\nu\) exactly when \(X=0\), \(\nu\)-almost surely.
The measures are mutually absolutely continuous exactly when
\(E_\nu X=c\) and \(X>0\), \(\nu\)-almost surely.
\end{hbtheorem}

\begin{hbtheorem}{thm:p15-partition-density-approximation}{proved=annex:p15.partition-density-approximation}{Partition ratios, a singular atom and insufficient information}
Suppose finite measurable partitions
\((I_{k,n})_{0\leq k<K_n}\) refine one another and generate
\(\mathcal F_n\). Then \(\mu_n\ll\nu_n\) exactly when
\(\nu(I_{k,n})=0\) implies \(\mu(I_{k,n})=0\).
A version of \(X_n\) on each cell is its mass ratio
\(\mu(I_{k,n})/\nu(I_{k,n})\), with value zero on null cells
\cite[Example~4.3.7]{DurrettPTEAuthor2019}.

For the dyadic partitions of \([0,1)\), Lebesgue probability \(\nu\),
\(x_0=1/3\), \(0<\theta<1\), and
\(\mu=(1-\theta)\nu+\theta\delta_{x_0}\), write \(I_n(x_0)\)
for the containing half-open dyadic cell of length \(2^{-n}\).
Then
\[
                 X_n=1-\theta+\theta2^n\mathbf1_{I_n(x_0)}.
\]
It tends to \(1-\theta\) at every \(x\ne x_0\), and to infinity
at \(x_0\). Its mean is one at every time and its limiting mean is
\(1-\theta\). Thus even a strictly positive finite limit
\(\nu\)-almost surely does not imply \(\mu\ll\nu\).

The generation hypothesis is necessary: with \(\mu=\delta_{x_0}\),
the same \(\nu\), and every \(\mathcal F_n=\{\varnothing,\Omega\}\),
the local derivatives are all one. They do not yield the Lebesgue
decomposition on the full Borel sigma-field.
\end{hbtheorem}

\section{Independent coordinates and the product criterion}

\begin{hbdefinition}{def:p15-product-density-laws}{Coordinate laws and square-root affinities}
Let \(\alpha_m,\beta_m\) be Borel probability measures on \(\mathbb R\),
\(m\geq1\), with \(\alpha_m\ll\beta_m\).
Let \(\mu=\bigotimes_{m\geq1}\alpha_m\) and
\(\nu=\bigotimes_{m\geq1}\beta_m\) on
\((\mathbb R^{\mathbb N},\mathcal B(\mathbb R)^{\otimes\mathbb N})\);
their existence and uniqueness are
\cref{thm:p15-independent-sequence-existence}.
Write \(\xi_m(\omega)=\omega_m\),
\(\mathcal F_n=\sigma(\xi_1,\ldots,\xi_n)\), and let
\(\mathcal F_0\) be the trivial sigma-field.
Choose finite nonnegative Borel densities \(q_m=d\alpha_m/d\beta_m\).
Put
\[
 h_m=\int_{\mathbb R}\sqrt{q_m}\,d\beta_m,\qquad
 H_n=\prod_{m=1}^nh_m,\quad H_0=1,\quad H_\infty=\lim_nH_n.
\]
The number \(h_m\) is the square-root affinity of the two coordinate
laws. All these integrals refer to the stated measures, not to
Lebesgue measure unless the latter is explicitly specified.
\index{Kakutani dichotomy}\index{affinity!square-root}
\end{hbdefinition}

\begin{hblemma}{thm:p15-product-density-filtration}{proved=annex:p15.product-density-filtration}{Finite likelihood products and affinities}
For \(n\geq0\),
\[
 X_n=\prod_{m=1}^nq_m(\xi_m)
       =\frac{d\mu_n}{d\nu_n},\qquad E_\nu X_n=1.
\]
The empty product \(X_0\) is one. For every \(m\),
\(0<h_m\leq1\), with equality \(h_m=1\) exactly when
\(\alpha_m=\beta_m\). Thus \(H_n\) decreases to a number in \([0,1]\).
For \(0\leq n<l\), putting \(V_n=\sqrt{X_n}\) gives
\[
 E_\nu V_n=H_n,\qquad E_\nu V_n^2=1,\qquad
 E_\nu(V_l-V_n)^2=2\left(1-\prod_{m=n+1}^lh_m\right).
\]
\end{hblemma}

\begin{hbtheorem}{thm:p15-product-positive-tail}{proved=annex:p15.product-positive-tail}{The source's strict-positivity tail dichotomy}
Suppose additionally \(q_m>0\), \(\beta_m\)-almost surely,
for every \(m\). Then the coordinate laws are equivalent and
versions may be chosen finite and strictly positive everywhere.
The reciprocal products \(R_n=1/X_n\) are restriction-density
martingales under \(\mu\). Both
\[
                  \nu(X=0)\in\{0,1\},\qquad
                  \mu(X=\infty)\in\{0,1\}
\]
follow from tail zero-one laws, where \(X=\limsup_nX_n\).
Consequently precisely one of \(\mu\sim\nu\) and
\(\mu\perp\nu\) holds.
The proof includes the reciprocal argument that determines the
singular mass; positivity of a limiting density alone is insufficient.
\end{hbtheorem}

\begin{hbtheorem}{thm:p15-product-affinity-criterion}{proved=annex:p15.product-affinity-criterion}{The square-root product decides absolute continuity}
Under \cref{def:p15-product-density-laws}, without requiring every
\(q_m\) strictly positive,
\[
 \begin{array}{lll}
 H_\infty>0 &\Longleftrightarrow& \mu\ll\nu,\\
 H_\infty=0 &\Longleftrightarrow& \mu\perp\nu.
 \end{array}
\]
In the first case \(X_n\to X=d\mu/d\nu\) in \(L^1(\nu)\).
In the second case \(X_n\to0\), \(\nu\)-almost surely, while
\(X_n\to\infty\), \(\mu\)-almost surely.

If all coordinate laws are equivalent, the first case gives
\(\mu\sim\nu\); otherwise its conclusion remains one-sided.
This retains the square-root proof of source Theorem~4.3.8 and
explicitly states the extension allowing zero coordinate densities
\cite[pp.~229--230]{DurrettPTEAuthor2019}.
\end{hbtheorem}

\begin{hbcorollary}{thm:p15-product-affinity-series}{proved=annex:p15.product-affinity-series}{An equivalent summability test}
With the same coordinate hypotheses,
\[
 \mu\ll\nu
 \quad\Longleftrightarrow\quad
 \sum_{m=1}^\infty(1-h_m)<\infty
 \quad\Longleftrightarrow\quad
 \sum_{m=1}^\infty
   \int(\sqrt{q_m}-1)^2\,d\beta_m<\infty.
\]
Every individual affinity is positive, so finite initial factors do
not change this criterion.
\end{hbcorollary}

\begin{hbtheorem}{thm:p15-product-gaussian-comparison}{proved=annex:p15.product-gaussian-comparison}{Full sequence laws for independent normal observations}
For \(\alpha_m=\operatorname{normal}(a_m,\sigma_m^2)\) and
\(\beta_m=\operatorname{normal}(b_m,\tau_m^2)\), with real means
and \(0<\sigma_m,\tau_m<\infty\), their affinity is
\[
 h_m=\sqrt{\frac{2\sigma_m\tau_m}{\sigma_m^2+\tau_m^2}}\,
 \exp\left(-\frac{(a_m-b_m)^2}{4(\sigma_m^2+\tau_m^2)}\right).
\]
The product laws are equivalent if \(\sum_m(1-h_m)<\infty\)
and singular otherwise. In particular when \(\sigma_m=\tau_m\),
\[
                \mu\sim\nu
       \quad\Longleftrightarrow\quad
           \sum_m\frac{(a_m-b_m)^2}{\sigma_m^2}<\infty.
\]
Equality of every finite-dimensional support does not settle
equivalence of the complete sequence laws.
\end{hbtheorem}

\begin{hbtheorem}{thm:p15-product-bernoulli-boundaries}{proved=annex:p15.product-bernoulli-boundaries}{Small coordinate changes and a zero-density boundary}
For Bernoulli laws with success probabilities \(p_m,q_m\in(0,1)\),
the affinity is
\[
                  h_m=\sqrt{p_mq_m}
                       +\sqrt{(1-p_m)(1-q_m)}.
\]
In this paragraph \(q_m\) is a scalar probability, not the earlier
density function. When \(q_m=1/2\) and
\(p_m=1/2+\varepsilon_m\), \(|\varepsilon_m|\leq1/4\),
\[
             \frac{\varepsilon_m^2}{3}
                  \leq1-h_m\leq\varepsilon_m^2.
\]
The laws are equivalent exactly when \(\sum_m\varepsilon_m^2<\infty\).
Thus \(\varepsilon_m=1/(4m)\) gives equivalence, whereas
\(\varepsilon_m=1/(4\sqrt m)\) gives singularity, although both
coordinate perturbations tend to zero.

If instead \(\nu\) is the fair Bernoulli sequence law and \(\mu\)
fixes only its first coordinate at zero, leaving the others independent
and fair, then
\[
 X_n=2\mathbf1_{\{\xi_1=0\}}\ (n\geq1),\qquad H_\infty=1/\sqrt2.
\]
Here \(\mu\ll\nu\), but \(\nu\not\ll\mu\) and
\(\nu(X=0)=1/2\). Strict coordinate positivity is necessary for
the tail assertion in \cref{thm:p15-product-positive-tail}.
\end{hbtheorem}

\begin{proofboundary}{def:p15-density-source-scope}{The exact presentation and its extensions}
The complete source chain here is Theorem~4.3.5, Lemma~4.3.6,
Example~4.3.7, the independent-product tail argument and Theorem~4.3.8.
All required density-sum, change-of-measure, chain and reciprocal
identities from Exercises~A.4.6--A.4.9 are supplied in
\cref{thm:p15-density-calculus}. The bounded mixture route above
preserves that source construction.

Allowing zero coordinate densities, the summability reformulation
and the normal and Bernoulli calculations are proved extensions.
The source's preceding general P\'olya beta limit is proved in
\cref{thm:p15-polya-beta-limit}. The distinct arguments of
Exercises~4.3.4--4.3.10, including the elementary density construction
and every Bernoulli product comparison, are supplied in
\cref{sec:p15-density-and-product-exercises}.
\end{proofboundary}
